\documentclass[a4paper,10pt]{article}
\usepackage[utf8]{inputenc}
\usepackage{graphicx}
\usepackage{subfigure}
\usepackage{fancyhdr}
\usepackage{listings}
\usepackage{ gensymb }


%opening

\pagestyle{fancy} %eigener Seitenstil
\fancyhf{} %alle Kopf- und Fußzeilenfelder bereinigen
\fancyhead[L]{SEP Abgabegespräch SS2015/16} %Kopfzeile links
\fancyhead[C]{} %zentrierte Kopfzeile
\fancyhead[R]{\includegraphics[width=1cm]{../../template/tugraz_bw}} %Kopfzeile rechts
\renewcommand{\headrulewidth}{0.4pt} %obere Trennlinie
\fancyfoot[C]{} %Seitennummer
\renewcommand{\footrulewidth}{0.4pt} %untere Trennlinie


\lstset{language=C++, basicstyle=\ttfamily\tiny}  

\begin{document}

\section*{Kostenmanager}

Ziel dieser Aufgabe ist es, eine Finanzverwaltungssoftware zu schreiben. Der User soll dabei Einnahmen und Ausgaben eintragen können und diese dann in einem angegebenen File mit der aktuellen Bilanz ausgeben. Folgende Befehle stehen zur Verfügung:
\begin{description}
 \item[quit]: Beendet das Programm.
 \item[add $<$typ$>$ $<$name$>$ $<$betrag$>$]: Fügt die Einnahme/Ausgabe mit einem Namen und einem Betrag ein.
 \item[save $<$file$>$]: Speichert die eingegebenen Daten in einem definierten Format.
\end{description}
Zeile für Zeile sollen Einnahmen und Ausgaben in eine Datei geschrieben werden.\newline
Ausgaben: \texttt{You spent <betrag> on <name>}.\newline
Einnahmen: \texttt{You received <betrag> from <name>}.\newline
In der letzten Zeile soll die aktuelle Bilanz (Summe aller Einnahmen und Ausgaben) abgespeichert werden.

\subsection*{Beispiel}
\footnotesize
\begin{verbatim}
./budget
> add income job 100
> add expense -30
Wrong parameter count!
> add expense groceries -30
> add expense rent -50
> save budget.txt
> quit
\end{verbatim}

Abgespeichertes File \texttt{budget.txt}:
\begin{lstlisting}[frame=single]
You received 100 from job
You spent -30 on groceries
You spent -50 on rent
Total: 20
\end{lstlisting}

\subsection*{Bewertung}
\begin{itemize}
 \item Dokumentation und Programmierstil: 1 Punkt
 \item Einhalten der OOP Konzepte: 1 Punkt
 \item Struktur \& Korrektheit: 1 Punkt
\end{itemize}
\newpage


\end{document}
